\chapter{材料・化学・バイオ}
\section{扱う対象と学習の出発点}
材料は組成だけでなく結合・組織・加工履歴で機能が変わる。化学は物質量・電荷・エネルギーの保存から変化を記述し、生命現象では分子が水中の反応系を構成する。まず何を数え、何に対する割合を取るかを固定すると、燃焼・濃度・核種・合金組成へ共通の見通しが得られる。以下は原子・化学量・機能性セラミックスと、分子から細胞の機能へつなぐ生命科学の初稿である。金属組織・物性・物質の性質も共通理論と全問の解説初稿へ接続する。

\section{原子・核種を数える}
\label{theory:basicchem:isotope}
原子番号\(Z\)は陽子数で元素を決め、質量数\(A\)は陽子と中性子の合計である。\({}^{A}_{Z}X\)の中性子数は\(N=A-Z\)、中性原子の電子数は\(Z\)、電荷\(+ke\)のイオンなら\(Z-k\)となる。質量数は核子数であり、実際の原子質量そのものではない。
\begin{equation}N=A-Z,\qquad n_{\mathrm e}=Z\quad\text{（中性原子）}.\end{equation}
同位体は同じ\(Z\)で異なる\(A\)を持つ核種である\autocite{iupac_isotopes}。Cl-35とCl-37は電子17個、中性子18個と20個である。Ca-40とAr-40は同じ\(A\)でも\(Z=20,18\)が違う同重体で、中性子20個と22個である。同素体は単体の結合・構造の違いで、核種の違いではない。
放射壊変で核の構成が変わると、例えば\(\beta^-\)壊変では中性子が陽子へ変わり\(Z\)が1増える。一方、励起核の\(\gamma\)遷移は\(Z\)を変えない。年代測定・トレーサー・診断等の利用を知るとともに、全壊変で別元素になるとは一般化しない。

\section{同じ電子配置と同位体効果}
\label{theory:basicchem:isotope_effect}
同位体の中性原子は核電荷と電子配置が同じなので、主要な結合様式・元素としての性質はほぼ共通である。しかし質量は分子振動へ入り、二原子振動子では
\begin{equation}\omega=\sqrt{k/\mu},\quad\mu=\frac{m_1m_2}{m_1+m_2},\quad E_0=\frac{\hbar\omega}{2}.\end{equation}
\(k\)は結合の力定数、\(\mu\)は換算質量である。質量変化は振動数・零点エネルギーを変え、速度や平衡定数の同位体効果を生む\autocite{iupac_isotope_effect}。試験の「化学的性質はほぼ同じ」という基礎的近似を、同位体効果が厳密に0という主張へ拡張しない。

\section{物質量は何を数える量か}
\label{theory:basicchem:amount}
反応係数は粒子個数の比を与え、実験では質量・気体体積を測る。物質量は両者をつなぐ。指定粒子の個数\(N_{\mathrm{ent}}\)、質量\(w\)、モル質量\(M\)について
\begin{equation}n=\frac{N_{\mathrm{ent}}}{N_{\mathrm A}}=\frac{w}{M},\quad N_{\mathrm A}=6.02214076\times10^{23}\ \mathrm{mol^{-1}}.\end{equation}
水なら分子、NaClなら式単位、合金なら各元素の原子と対象を明記する。\(N\)を核種の中性子数と混同しない。理想気体では\(n=pV/(RT)\)。0℃、約1気圧なら約22.4 L/molだが、別の温度・圧力へそのモル体積を持ち込まない。現行の定数定義と歴史的な測定精度を分け、指定された原子量・丸め値を計算へ用いる。

\section{反応係数と保存則から収率を求める}
\label{theory:basicchem:stoichiometry}
化学反応では通常の原子核は変わらず、元素ごとの原子数が保存される。\(\Delta n_i=\nu_i\xi\)では\(\nu_i\)が生成側を正、反応側を負に取る反応係数、\(\xi\)が反応進行度である。分子の炭素数が\(c\)、完全燃焼で全炭素がCO2になるなら
\begin{equation}n_{\mathrm{CO_2}}=c\frac{w}{M}.\end{equation}
同じ質量\(w\)の比較は\(c/M\)の比較となる。炭素数だけでなく水素・酸素の質量も分母へ入る。部分燃焼、すす生成、酸素不足は前提の外である。
酸素消費量は生成物に依存する。\(2\mathrm{Cu}+\mathrm{O_2}\to2\mathrm{CuO}\)と\(4\mathrm{Cu}+\mathrm{O_2}\to2\mathrm{Cu_2O}\)では必要量が2倍違う。生成物が未指定なら仮定を示し、可能な生成物でも順位が変わるか確認する。

\section{濃度の分母を先に決める}
\label{theory:basicchem:concentration}
濃度には何で割るかを含む。溶質の物質量\(n_B\)、最終溶液体積\(V\)、溶媒質量\(w_{\mathrm{solvent}}\)について
\begin{equation}c_B=\frac{n_B}{V}\quad[\mathrm{mol/L}],\qquad b_B=\frac{n_B}{w_{\mathrm{solvent}}}\quad[\mathrm{mol/kg}].\end{equation}
前者はモル濃度、後者は質量モル濃度である\autocite{iupac_amount_concentration,iupac_molality}。水1 Lへ溶質1 molを加えても最終溶液が1 Lとは限らない。1 mol/Lの調製は溶解後に全体を1 Lへ定容する。体積は温度・混合で変わるため、溶媒kgを分母とする量とは異なる。
質量10\%は溶質100 gと溶媒900 gから1000 gの溶液を作る場合である。質量ppmは質量分率の\(10^6\)倍。1 mg/Lは質量濃度で、その定義がppmではない。密度約1 kg/Lの希薄水溶液なら質量約1 ppmに一致する。近似と密度・基準の省略を無条件の等号へ変えない。

\section{質量分率とモル分率の変換}
\label{theory:basicchem:fraction}
モル分率は全成分の物質量に対する成分\(i\)の割合である\autocite{iupac_amount_fraction}。
\begin{equation}x_i=\frac{n_i}{\sum_jn_j}=\frac{w_i/M_i}{\sum_j w_j/M_j},\qquad\sum_i x_i=1.\end{equation}
合計100 gを仮に取ると、質量百分率をそのままgとして計算できる。仮の全体質量は分子・分母で消える。原子量が小さい成分は同じ質量でより多くの原子を持つので、質量分率からモル分率へ変換したとき増える方向が期待できる。合金では原子、溶液では分子・式単位と、対象を指定する。

\section{酸化数を電荷収支から求める}
\label{theory:basicchem:oxidation}
酸化数は異種原子間の結合をイオン的に近似した電子配分の指標であり、実際の電荷密度そのものではない\autocite{iupac_oxidation_state}。今回の酸・酸化物はH=+1、O=−2とし、総和を全体電荷へ等置する。
\begin{equation}\sum_i a_i s_i=q.\end{equation}
\(a_i\)は原子数、\(s_i\)は原子1個の酸化数、\(q\)は電気素量単位の電荷である。MnO4\({}^-\)なら\(x-8=-1\)から+7、Cr2O7\({}^{2-}\)なら\(2x-14=-2\)からCrは+6である。単体Mnは0、NH3のNは−3となる。2個のCrの合計+12と1原子の+6を分ける。
過酸化物のO=−1、金属水素化物のH=−1等の例外を今回の化合物が含まないことを確認して使う。酸化数増加が酸化、減少が還元であり、酸素が式にあるかだけでは判定しない。

\section{多原子イオンと式単位}
\label{theory:basicchem:ion}
イオンの価数と個数は違う。NH4\({}^+\)、OH\({}^-\)、NO3\({}^-\)は多原子1価、SO4\({}^{2-}\)は多原子2価、PO4\({}^{3-}\)は多原子3価である。S\({}^{2-}\)、Ca\({}^{2+}\)、Pb\({}^{2+}\)は2価でも単原子である。
Pb(NO3)2の括弧外2はNO3\({}^-\)が2個あることを表し、硝酸イオン1個を−2へ変えない。(NH4)2SO4は2個の+1と−2、Ca3(PO4)2は3個の+2と2個の−3が釣り合う。全体の中性から確かめると名称暗記に頼らず判断できる。

\section{結合エネルギーと反応エンタルピー}
\label{theory:basicchem:enthalpy}
結合切断にはエネルギーを与え、生成では放出する。反応物を気相原子へ分け、生成物の結合を作る仮想経路をヘスの法則で合計する\autocite{mit_bond_enthalpy}。
\begin{equation}\Delta H\simeq\sum D_{\mathrm{broken}}-\sum D_{\mathrm{formed}}.\end{equation}
\(D\)は正の切断エンタルピーである。平均結合エンタルピーは化学環境の違いによる近似を含み、液体・固体の反応では相変化等も必要となる。H2+Cl2\(\to\)2HClで\(D_{HH}=436\)、\(D_{ClCl}=243\)、\(D_{HCl}=432\) kJ/molなら\(\Delta H=679-864=-185\) kJ。HClを2 mol生成する基準なので1 mol当たり−92.5 kJである。
発熱は系の\(\Delta H<0\)だが「外へ発生する熱の大きさ」は正で表すこともある。原典の符号指定と候補を読み、平成23年度の放出熱量と平成26年度の発熱負の指定を区別する。

\section{ハロゲンを一つの順位で表さない}
\label{theory:basicchem:halogen}
電気陰性度は結合中の電子を引く傾向、酸化力は電子を受け取る反応の駆動力、酸性度はプロトン移動の平衡、沸点は凝集相から気相へ移る条件である。同じ元素列でも別現象である。電気陰性度と分子の酸化力は概ねF、Cl、Br、Iの順に減少する\autocite{purdue_halogens,rsc_halogen}。
しかし水中のHFは弱酸、HCl・HBr・HIは強酸である。結合切断の費用、陰イオンの安定性、溶媒和を含むので電気陰性度だけではHXの酸性度を決められない\autocite{purdue_acid_strength,yale_hydrohalic}。同じ分析濃度の希薄水溶液でHFのH\({}^+\)が最大という肢は誤り。強酸3種はほぼ完全解離し、水による水平化で単純な濃度差としては区別しにくい。
通常沸点はHF\(>\)HI\(>\)HBr\(>\)HClである\autocite{nist_halogen_properties,msu_hydrogen_bonding}。HFは水素結合のため特に高く、他3種は大きい分子ほど分散力が強くなる。X2の沸点とHXの沸点も別々に読む。

\section{電子セラミックスを入出力から理解する}
\label{theory:basicchem:ceramics}
機能材料では、何を入力すると何が変わるかを対応させる。BaTiO3系誘電体は高い比誘電率により容量を大きくする\autocite{murata_batio3}。\(C=\varepsilon_0\varepsilon_r A/d\)は幾何学と分極応答を分ける。導電率を上げて電荷を通すことは電極間を絶縁して蓄える機能と異なる。
サーミスタは温度で抵抗が変わる素子である\autocite{murata_thermistor}。圧電体は応力・ひずみと電気応答を結び、変形に伴う電荷を電圧へ変換してセンサーとする\autocite{murata_piezo}。熱膨張とは異なる結合項を使う。
バリスタは非線形な電圧–電流特性を使う\autocite{tdk_varistor}。\(I=KV^\alpha\)の近似なら\(V/I=K^{-1}V^{1-\alpha}\)で、\(\alpha>1\)では電圧上昇で実効抵抗が低下する。\(\alpha=1\)の純抵抗と区別し、サージ時に電流を逃がす機能へ利用する。温度抵抗・圧電・誘電・非線形抵抗を別々の入力出力として理解する。

\section{生命現象を組成・区画・情報から見る}
\label{theory:biology:composition}
生命科学では、同じ化学結合を使って反応を選択し、情報を保存・複製し、細胞の境界を維持する仕組みを扱う。細胞は水中の反応系である。糖・脂肪酸・アミノ酸・ヌクレオチドという小分子の群を、エネルギー源、膜材料、高分子の単位として区別すると用語に役割が付く。
典型的な細胞で水が総重量の約70\%という概数を用いる。一方、H・O・C・Nの主要元素やHの約60\%以上という説明は原子の個数を分母とする。水だけなら3個の原子のうち2個がHであるが、質量ではHは約2/18にすぎない。軽い元素の個数割合と質量割合は一致しない。
\begin{equation}
 x_i^{\rm atom}=\frac{N_i}{\sum_jN_j},\qquad
 w_i=\frac{N_i M_i}{\sum_jN_jM_j}.
\end{equation}
ここで\(N_i\)は元素\(i\)の原子数、\(M_i\)はそのモル質量である。同じ比例定数で原子数を物質量へ換算できるため、式の比では定数が消える。
通常の細胞の有機成分・巨大分子の重量ではタンパク質の寄与が大きい。実験研究が比較する細菌・酵母・哺乳類細胞の組成も同一ではなく、培養状態で変わる\autocite{bio_cell_mass}。典型例の数字は全種・全組織の固定比率ではない。平成26年度問4-6の「すべての生物で同様」はこの範囲に疑義があり、公式の主要元素の誤りと別に追跡する。

\section{膜脂質：水を避ける鎖が境界を作る}
\label{theory:biology:membrane}
リン脂質は親水的な頭部と疎水的な炭化水素鎖を併せ持つ。水中で頭部を水へ向け、鎖を内側に集めることで二重層を作る。通常のグリセロリン脂質では脂肪酸がエステル結合でグリセロールへつながる。脂肪酸の全結合形式がエステルという定義ではなく、スフィンゴ脂質のアミド結合、古細菌のエーテル脂質を区別する。
膜は完全に固定した壁ではなく、脂質等が面内を移動する境界である。温度、鎖長、不飽和度が充填と運動性を変える。同系統の飽和鎖では長い鎖ほど相互作用が増し、融点が高くなる。cis二重結合は鎖を折り曲げ、密な充填を妨げ、同炭素数の飽和鎖より流動性を増す方向に働く\autocite{bio_membrane}。trans型や混合物の全挙動をこの一つの近似へ押し込めない。
温度上昇そのものは流動化する。環境温度が変わる細菌は長鎖化・飽和化等の組成変更でこれを相殺する場合がある。設問の「増える場合がある」を、全細菌で同じ割合になるという法則へ変えない。膜の適切な流動性は輸送・酵素機能を維持する条件であり、単に流動性を最大にすればよいわけではない。

\section{アミノ酸：共通骨格と側鎖を分ける}
\label{theory:biology:amino}
タンパク質を形作る標準20種のアミノ酸は、α炭素にアミノ基、カルボキシ基、H、種類ごとのR側鎖を持つ。水中の電荷状態はpHに依存し、中性付近ではアミノ基の正電荷とカルボキシ基の負電荷を併せ持つ双性イオンとなる場合が多い\autocite{mit_amino}。化学式に未解離のNH2・COOHを書くことと、溶液内の主な状態を区別する。
ロイシン・イソロイシン等の炭化水素側鎖、フェニルアラニンの芳香族側鎖は非極性で疎水的である。システインのRはチオールを持ち、メチオニンのRはチオエーテルを持つ。両者に硫黄があることと、両者が同じS-S架橋を作ることは別である。
グリシンはR=Hでα炭素が不斉でない。他の標準タンパク質アミノ酸は主にLである。D/Lは立体配置の系列を表し、旋光の正負を表す記号ではない。「特殊なものを除く」という出題条件を確認すれば、標準20種と修飾・特殊アミノ酸を混ぜず判断できる。

\section{タンパク質：配列から構造と機能へ}
\label{theory:biology:protein}
隣接するアミノ酸のカルボキシ基とアミノ基に由来する共有のペプチド結合が主鎖を作る。配列を一次構造、局所的なαらせん・βシート等を二次構造、全体の折畳みを三次構造、複数鎖の集合を四次構造という\autocite{mit_protein}。階層の名前と、その維持に働く相互作用の名前を分ける。
配列は翻訳する核酸配列に規定される。一方、機能する立体構造には溶液条件、シャペロン、結合相手等も関わるので、配列の指定を全環境で同じ活性が保証されるという意味にしない。主鎖のペプチドとCys間のジスルフィドは共有結合である。静電相互作用、水素結合、分散力等の非共有相互作用も折畳みへ寄与する。
疎水効果は水中で非極性面を露出する自由エネルギー費用に関係する。通常の水溶性球状タンパク質では非極性側鎖が内部へ偏るが、膜の脂質に接する面などは違う。疎水性相互作用を新しい主鎖の共有結合とはしない。誤答肢の「ジスルフィド等の非共有結合」は、構造への寄与が正しくても結合分類が誤りである。

\section{等電点：正味0と各基の無電荷は違う}
\label{theory:biology:charge}
タンパク質は末端と側鎖に複数の解離基を持ち、それぞれのプロトン化率がpHで変わる。正味電荷0となるpHが等電点pIである。通常は\(pH<pI\)で正、\(pH>pI\)で負となる。酸塩基解離でpH上昇に伴いプロトンを失う方向を追うと符号を理解できる。電気泳動や分離では溶液のpHとpIを比較する\autocite{mit_protein}。
正味0は正負の電荷が相殺することであり、全ての官能基が無電荷という条件ではない。多くの解離基を持つタンパク質のpIを、二つのpKaの平均だけで常に求められるとはしない。

\section{DNA：鎖の内部と鎖同士の会合}
\label{theory:biology:dna}
ヌクレオチドは塩基・糖・リン酸からなる。DNAの糖はデオキシリボース、標準塩基はA・T・G・Cである。RNAはリボースと通常Uを用いる。各鎖でヌクレオチドをつなぐ糖リン酸主鎖はホスホジエステル共有結合であり、二重らせんで塩基は内側、主鎖は外側となる。
A-TとG-Cが相補的に会合し、2鎖は逆平行である\autocite{nhgri_dna}。水素結合は選択性へ、塩基スタッキング等も安定化へ寄与する。会合が解けることと主鎖が分解することを別に扱う。
通常のDNAポリメラーゼは鋳型とプライマーの自由な3端を使ってヌクレオチドを付ける。したがって新生鎖は\(5'\to3'\)、鋳型の読みは逆方向となる。「DNAが二重鎖」の典型説明を、一本鎖DNAを持つウイルス等が存在しないという主張にはしない。

\section{塩基組成：一方の鎖に全体の対称性を課さない}
\label{theory:biology:base_fraction}
一方の鎖の割合を\(a,t,g,c\)、相補鎖を星付きとすれば
\begin{equation}
 a+t+g+c=1,\qquad
 a^{*}=t,\ t^{*}=a,\ g^{*}=c,\ c^{*}=g.
\end{equation}
相補鎖は各塩基の行先を入れ替える写像である。一方の鎖だけに\(a=t\)や\(g=c\)を課す根拠はない。二重鎖全体では対を合算するためその等式が成り立つ。
C20\%、T35\%なら同じ鎖のA+G45\%、相補鎖A+G55\%である。G25\%、A15\%なら相補鎖C25\%、T15\%、G+A60\%である。個々の未指定値は確定しないが、非負と総和1の制約で不可能な肢は排除できる。例えば後者の同じ鎖G+C90\%はC65\%を要求し、C+T60\%に反する。

\section{DNA融解と吸光度：配列を切らずに会合を解く}
\label{theory:biology:melting}
熱又は適切な化学条件で二重鎖は一本鎖へ解離する。通常の変性で主鎖のホスホジエステル結合は残る。変性に伴い約260 nmの紫外吸収が増す増色効果を用いて融解曲線を観測する\autocite{bu_dna_melting}。二重鎖の塩基スタッキングの変化が吸光度へ影響する。
転移の中点を融解温度\(T_m\)という。二状態近似であれば半分が解離した点に対応する。長い不均一なDNAは領域ごとに転移し得るので、分子全体の個数が常に厳密に半数という説明へ固定しない。他条件を揃えるとGC含量が高いほど一般に\(T_m\)は高いが、塩濃度、鎖長、配列、ミスマッチ等も影響する。
相補鎖が存在し、濃度・温度・時間が適切なら冷却で再会合する。これがアニーリングである。物質の共有結合分解を逆に作り直す操作とは異なる。この区別がPCRの熱変性とプライマー会合の理解へ直接つながる。

\section{遺伝暗号：多対一の対応と三文字の読み枠}
\label{theory:biology:code}
転写はDNAの情報からRNAを作り、翻訳ではmRNAを3塩基ずつ読んでアミノ酸鎖を作る。4文字の3文字列は\(4^3=64\)通りであるが、入力の個数と出力種数は一対一である必要はない。標準暗号では61が20種のアミノ酸、3が終止に対応する\autocite{nhgri_codon}。複数のコドンが同じアミノ酸を指定することを縮重という。
残る44文字列が消失した、44の修飾アミノ酸を指定する、転写で20へ書き換わるという説明は違う。縮重があることから第1塩基が無意味とは導けない。第3位置の置換で同義となる場合も全てではない。開始位置が読み枠を定め、同じ文字列でも区切りが変われば指定配列が変わる。
この標準表は特殊な再符号化やミトコンドリア等の暗号差を別にして用いる。DNAのコドンという設問表記は、転写後のmRNAでTに対応してUが使われることと区別する。

\section{突然変異：文字を替えることと区切りをずらすこと}
\label{theory:biology:mutation}
塩基置換は一文字を替え、挿入・欠失は文字数を変える。通常の翻訳領域で3の倍数でない文字数の変化は下流の読み枠をずらすフレームシフトとなる\autocite{nhgri_frameshift}。非翻訳領域の全挿入を、同じ読み枠の式だけで扱わない。
置換は同義となる場合、アミノ酸が変わるミスセンス、途中の終止となるナンセンス等を生む。代表的なHbSのGlu→Valは置換であって1塩基欠失ではない\autocite{nhgri_sickle}。中立性は機能や選択への効果を論じる軸で、ナンセンスと同義の名前ではない。
常染色体の通常の二倍体・完全潜性のモデルでは、正常な対立遺伝子が一方にあると形質が隠れる。両方の状態を見る必要がある。半接合、不完全な優性、組織モザイク等は前提が違う。X線・UVや特定の化学物質はDNA損傷・変異の誘発源となり得るが、物質群の全てを一律に変異原としない。

\section{PCR：温度で分け、目印から増やす}
\label{theory:biology:pcr}
目的領域の両端へ対応する2本のプライマーを用い、熱変性、アニーリング、伸長を繰り返す。DNA主鎖を切るのではなく鎖間会合を解き、プライマーの3端から合成する。プライマー自身の配列は産物の一部となる。
反復加熱に耐えるDNAポリメラーゼを用いる。アニーリング温度が低すぎると非特異的会合が増え、高すぎると目的会合も減り得る。実際の温度・時間はプライマー配列、塩条件、酵素製品等に応じて設定する\autocite{neb_pcr}。伸長速度を\(u\)塩基/秒、標的長を\(L\)塩基とすれば、最小限の時間は概ね\(L/u\)であり、長い配列ほど短くするという判断は逆である。
完全な複製を各回仮定する理想モデルでは
\begin{equation}
 N_{n+1}=2N_n,\qquad N_n=N_0\,2^n.
\end{equation}
\(N_n\)は\(n\)サイクル後の二重鎖相当の分子数で、30回は約\(1.07\times10^9\)倍となる。正確な両端長の産物が最初のサイクルから全分子を占めるわけではなく、実反応は効率・基質枯渇・阻害等でこの理想から外れる。回数をそのまま増幅倍率へ読み替えない。

\section{DNAライブラリと制限酵素：保存された情報と発現した情報}
\label{theory:biology:libraries}
ゲノムライブラリはゲノムDNAの断片群を、cDNAライブラリは採取したRNAから逆転写で作るDNA群を保持する。cDNAは発現したRNAに由来する\autocite{nhgri_cdna}。ゲノムにある遺伝子でも、その組織・時点で未発現なら対応するcDNAはない場合がある。成熟mRNA由来のcDNAでは通常イントロンが除かれるが、UTRは含み得るので「翻訳領域のみ」とは定義しない。
通常の緩衝液でDNAのリン酸主鎖は負に帯電し、ゲル電気泳動では陽極へ進む。分離は長さ等に依存し、陰極へという肢は電荷と移動方向を逆にしている。
6塩基の特定認識配列について、各位置が独立で4塩基が等確率という仮想モデルでは出現確率は\(4^{-6}=1/4096\)である。これは認識開始位置の平均間隔のモデルで、制限酵素切断後の全断片が正確に4096塩基対という主張ではない。ゲノムの組成偏り、配列相関、末端、切断可能性も別に考える。EcoRIの製造者文書も個別の酵素条件を扱う\autocite{neb_ecori}。
組換えインスリンの製法開発は1978年、米国での製品承認は1982年である\autocite{fda_insulin}。実験の成功・製品承認・後年の類似薬を同じ時点にせず、1990年代後半という誤答を識別する。

\section{酵素：平衡を変えず反応経路を変える}
\label{theory:biology:enzymes}
酵素は基質に選択的に作用する生体触媒である。多くはタンパク質だが触媒RNAもある。触媒は低い障壁の反応経路を与え、速度を変える。反応物と生成物の自由エネルギー差や平衡定数を変えることとは異なる。単純な活性化過程の近似\(k\propto\exp[-\Delta G^\ddagger/(RT)]\)では、障壁低下が速度増加へつながる。
リパーゼはトリアシルグリセロールのエステル結合を加水分解する酵素、αアミラーゼはデンプン等のα1→4結合へ作用する酵素である\autocite{iubmb_lipase,iubmb_amylase}。セルロースのβ1→4結合は同じグルコースからなるという理由だけでアミラーゼの標的とはならない。単位分子だけでなく結合の型が酵素の選択性を決める。
温度を上げると反応が速くなる効果と、酵素が変性・失活する効果が競合するため、至適温度と耐熱性は別概念である。Pyrococcus woeseiのαアミラーゼには100℃至適の一次研究があり、20℃以下又は100℃以上の酵素は一切ないという肢を反証する\autocite{hyper_amylase}。一つの存在例で不存在の全称命題を否定できる。

\section{クローン：分化状態を変え、核情報を受け渡す}
\label{theory:biology:cloning}
植物の組織培養ではオーキシン・サイトカイニン等のホルモン条件を用いてカルス形成や器官の再分化を制御する。ポプラの実験研究もホルモン処理とカルス・シュートの段階を比較している\autocite{plant_regeneration}。カルスは設問では不定形の増殖細胞塊、プロトプラストは細胞壁を取り除いた細胞であり別の語である。全カルスの全細胞が完全に同一の未分化状態という一般定義へ拡張しない。
哺乳類の体細胞核移植では、核を除いた卵へ体細胞核を移し、適切な操作で発生を進める\autocite{nhgri_cloning}。核を胚と読み違えたり受容側を卵原細胞とすると手順が成立しない。卵細胞の環境が核の発生プログラムを変えることが必要である。
体細胞クローンの核DNAの由来と、卵由来のミトコンドリアDNA、エピゲノム、発生環境を区別する。核情報が共通であっても外見・機能・健康状態が全て同一とは限らない。

\section{メチル化：配列を替えず使用状態を伝える}
\label{theory:biology:epigenetics}
同じゲノムを持つ細胞でも発現する遺伝子の組合せが違い、異なる役割を担う。DNAメチル化やヒストン修飾は配列を変更せず、その使用状態へ関わる。メチル化状態の維持は細胞分裂後の発現パターンの記憶へ寄与する\autocite{nhgri_epigenomics}。
これは一塩基置換のような配列変異と異なる。全ての修飾が必ず維持される、全遺伝子のメチル化が同じ効果を持つ、個体の全特徴が環境から子孫へ継承されるという主張ではない。設問の「子孫細胞」は細胞分裂での継承として読み、親子個体の世代を無条件に追加しない。

\section{必須アミノ酸：重要性でなく供給条件}
\label{theory:biology:nutrition}
必須アミノ酸は、生体が必要量を体内合成で満たせず外部供給を要するものをいう。非必須もタンパク質の機能へ重要であり、重要でないという分類ではない。通常のヒトの必須9種はHis・Ile・Leu・Lys・Met・Phe・Thr・Trp・Valである。Proは通常非必須に分類され、病態・成長等の条件では条件付き必須となり得る\autocite{nih_amino}。
令和7年度のIle・Val・Leu・Pro・Metでは通常の分類によりProを選ぶ。設問の広い「哺乳類」という語を、全種・全発育段階で栄養条件が同じとする保証へ変えない。前駆体から作れることと、その状況で必要量を十分作れることも区別する。

\section{出題形式と確認手順}
今回の18問では粒子数、同位体の正誤表、酸化数の最大・最小、燃焼収率、濃度調製、組成変換、結合エネルギー、材料語句の組合せを確認した。全年度の頻度は残る分類確認後に更新する。対象と分母を決め、保存則で計算し、候補順と否定方向を確認する。正答番号が合っても密度近似・生成物未指定の条件不足は残す。

\section{過去問}
\section{金属材料：結合・結晶・組織から試験と選択へ}
\label{theory:material:overview}
材料選択では、単に「強い金属」を探すだけでは足りない。小さな荷重での変形、永久変形の始まり、き裂の進展、繰返し負荷、電気抵抗、加工のしやすさは異なる特性である。組成と加工条件が結晶・転位・粒界・欠陥を変え、それが使用中の応答へ結び付く。本節は基礎科目の材料問題から、その因果関係を整理する。

\subsection{金属結合と塑性：電子と転位の役割}
\label{theory:material:bonding}
金属の外側の電子は原子一つに局在せず、正の原子核・内殻電子からなる芯の間で共有される。自由電子の海という近似は金属結合と導電性を理解する入口である\cite{mit_material_atomistic2001}。同位体は中性子数の違う同じ元素の原子であり、自由電子の代わりの結合機構ではない。

結合が特定方向だけに限られないことと、原子面が容易に全て同時に滑ることは同じではない。実際の塑性では格子の線欠陥である転位が移動し、少しずつすべりを進める。転位の移動を妨げる粒界・溶質・析出物は強化に寄与するが、延性との交換条件も生む\cite{mit_material_dislocations2001}。結合だけで全金属の室温延性を保証しない。

\subsection{BCC・FCC・HCP：胞と配位数を数え分ける}
\label{theory:material:crystals}
単位胞は結晶の並進で繰り返す領域である。格子点と、その各点に付く原子の基底を区別する。胞の形を決めずに「一胞に何原子」と尋ねると、原始胞と大きな慣用胞の数が混ざる\cite{cambridge_unitcell,princeton_hcp}。

通常の立方胞では、角の原子は隣り合う8胞で共有し、面の原子は2胞で共有する。この分率から独自に数えると、
\[
 N_{\rm BCC}=8\times\frac18+1=2,\qquad
 N_{\rm FCC}=8\times\frac18+6\times\frac12=4.
\]
BCCの最近接配位数は8、FCCは12である。配位数は一つの原子の周囲にある最近接原子の数であり、胞の中の原子数ではない\cite{openstax_crystals}。

硬球近似で半径を\(R\)、立方胞辺を\(a\)とする。BCCは体対角線、FCCは面対角線で接触するため、幾何から
\[
 \sqrt3a=4R\quad({\rm BCC}),\qquad
 \sqrt2a=4R\quad({\rm FCC})
\]
となる。FCCでは\(a=2\sqrt2R\)である。式の暗記だけでなく、接触する原子列を指定する。

HCPは六方Bravais格子に二原子基底を置いた構造であり、原始並進胞に2原子を含む。よく図示される大きな六角柱では、角12原子の共有分\(12/6\)、上下の面中心2原子の共有分\(2/2\)、内部3原子を合計し、6原子となる。この六角柱は原始並進胞の3倍の体積を持つ。平成29年度と令和6年度の「HCPは2原子」は前者の規約で読み、六角柱の図へ無条件に転記しない。

常温常圧の元素ではFe・NaはBCC、Al・CuはFCC、Ti・ZnはHCPである\cite{princeton_iron,princeton_sodium,princeton_aluminum,princeton_copper,princeton_titanium,princeton_zinc}。相の安定性は温度・圧力・組成に依存する。別資料の温度を指定しない元素例示を、室温の分類にそのまま使わない。

\subsection{荷重と変位から応力・ひずみへ}
\label{theory:material:tension}
初期断面積を\(A_0\)、初期標点距離を\(L_0\)、荷重を\(F\)、標点間の長さを\(L\)とする。公称値は
\[
 \sigma_{\rm e}=\frac{F}{A_0},\qquad
 \varepsilon_{\rm e}=\frac{L-L_0}{L_0}
\]
である。試験機のクロスヘッド変位は機械・つかみの変形を含むことがあるため、試験片の標点間の伸びと対応させる。

現在の断面積\(A\)で割った\(\sigma_{\rm t}=F/A\)が真応力である。真ひずみは現在長さによる微小増分を積分するので、
\[
 \varepsilon_{\rm t}=\int_{L_0}^{L}\frac{d\ell}{\ell}
 =\ln\frac{L}{L_0}=\ln(1+\varepsilon_{\rm e}).
\]
全伸びを変形後の長さで割るだけでは、この積分値にならない\cite{mit_material_stressstrain2001}。微小ひずみでは両ひずみは近い。大変形・くびれ後には、標点間平均と局所値を区別する必要がある。

\subsection{フック則：剛性と強度を分ける}
\label{theory:material:elasticity}
小応力の単軸線形弾性では\(\sigma=E\varepsilon\)となる。比例定数\(E\)がヤング率であり、ひずみが無次元なので応力と同じPaを単位に持つ。棒の剛性は\(k=EA_0/L_0\)であり、材料の\(E\)に加えて形状にも依存する\cite{mit_material_elastic2000}。

降伏応力は塑性が始まる応力、引張強度は通常、公称応力の最大値である。ヤング率はこの二つの強度ではない。ポアソン比は横と軸のひずみの比であり無次元である。ヘンリー則は溶解、ブラッグ則は回折の関係である。

出題のMg45 GPa、W407 GPaは大学教材にもある代表値である\cite{princeton_magnesium,princeton_tungsten}。同じ応力ならMgの方が大きく弾性変形する。通常の金属では昇温に伴う弾性率低下を考えるが、方位・温度区間・相変態を指定し、全材料の単調減少へ一般化しない。弾性は除荷による回復を指し、線形性と常に同義ではない。

\subsection{粒径・加工硬化・再結晶}
\label{theory:material:strengthening}
一般的な多結晶鋼などでは、結晶粒界が転位移動の障害になる。通常の粒径範囲の経験的関係は
\[
 \sigma_y=\sigma_0+\frac{k_y}{\sqrt d}
\]
である。\(d\)が小さいほど第二項が大きくなる。粗大化で強度が上がるという出題肢は、この通常の比較では逆である。

塑性加工で転位が増えると、転位同士も移動を妨げる。これが加工硬化を理解する基本である。加熱による回復・再結晶は蓄積した欠陥を減らし、通常は加工硬化を緩和して延性を回復する\cite{mit_material_dislocations2001}。

Hall--Petch式を原子数が極端に少ないナノ粒まで外挿し、強度が際限なく増すとはしない。粒界すべり・回転・拡散や加工時の欠陥など、支配機構が変わる。逆Hall--Petchの報告にも加工・硬度から降伏応力への換算・試験中粒成長の問題があり、全てのナノ材料が必ず弱化するという保証にしない\cite{hall_petch_limits2020}。

\subsection{欠陥・脆性破壊・疲労・衝撃靭性}
\label{theory:material:fracture}
無欠陥の原子結合を一斉に切る理論強度と、孔・介在物・き裂を含む実材料の強度は異なる。き裂先端には応力が集中し、き裂寸法・形状・負荷と破壊抵抗を一緒に評価する。線形破壊力学の典型的形は\(K_I=Y\sigma\sqrt{\pi a}\)であり、\(a\)の定義と形状係数\(Y\)を指定する。

脆性破壊は顕著な塑性変形が小さい破壊、疲労は繰返し負荷による損傷蓄積である。疲労では一回なら耐えられる荷重でも、き裂が成長して最後に急速破断することがある\cite{mit_material_fatigue2001}。腐食疲労はあるが、疲労を経年腐食だけで定義しない。

リバティ船とコメットは破壊力学発展を促した代表例である\cite{nist_fracture_history2010}。出題は船の脆性と航空機の疲労を対照する。コメットに関する大学教材は窓周辺のリベット孔からのき裂も説明する\cite{mit_material_fracture2001}。原典の「窓の角」という簡略表現を全事故の起点の保証へ変えず、溶接部・開口形状・低温・繰返し圧力などの複合条件を考える。

Charpyは切欠き試験片を衝撃破断し、吸収エネルギーから衝撃靭性を評価する。温度・寸法・切欠き・負荷速度に依存する測定であり、平面ひずみ破壊靭性\(K_{IC}\)の直接測定ではない\cite{nist_charpy}。

\subsection{純鉄の結晶相と磁気転移}
\label{theory:material:ironphases}
常温の純鉄\(\alpha\)-FeはBCCで強磁性、Coも常温の強磁性金属である\cite{princeton_iron,openstax_magnetism}。加熱時には約770°Cで強磁性を失い、BCCのまま常磁性となる。Curie温度1043 Kを摂氏へ換算すると約770°Cである\cite{iron_curie1954}。

その後、約912°Cで\(\gamma\)-FeのFCCへ変態する。加熱時の\(\alpha\to\gamma\)は収縮し、冷却時の逆変態では膨張する\cite{iron_transform2023}。磁気転移と結晶の同素変態は別であり、強磁性が結晶変態温度まで続くとしない。一般の熱膨張と変態による急な変化も別に考える。原子数と実際の格子定数が体積を決めるので、充填率だけで全材料の変態体積を決めない。

\subsection{被膜質量：寸法と百分率の分母}
\label{theory:material:coating}
面積\(S\)、被膜厚\(t_c\)、基材厚\(t_s\)に対し、片面被膜の質量は\(m_c=St_c\rho_c\)、基材は\(m_s=St_s\rho_s\)である。被膜の質量百分率は\(100m_c/(m_c+m_s)\)となる。厚さを同じ長さ単位にそろえ、片面・両面を図示条件に合わせる。

令和4年度基礎4-3の原典値を独自計算すると、
\[
 m_{\rm Ni}=0.50\times1.0\times10\times10^{-6}\times8900
 =0.0445\ {\rm kg},\qquad
 m_s=0.50\times1.0\times0.60\times10^{-3}\times7900
 =2.37\ {\rm kg}.
\]
したがって\(100\times0.0445/(2.37+0.0445)=1.843\ldots\%\)である。候補に対応する丸め値は\(4.5\times10^{-2}\)kg、1.8\%となる。百分率は鋼だけを分母にしない。

Niのレアメタル分類は日本の政策文脈であり、JOGMECは国際的にはベースメタルと呼ばれることも明記する\cite{jogmec_nickel2014}。分類語を普遍的な自然科学上の属性と思わず、計算条件とは別に確認する。

\subsection{材料選択：軽さ・剛性・加工性・使用環境}
\label{theory:material:selection}
軽量化でも、破断・永久変形を避ける選択と、弾性変形を抑える選択は異なる。比強度は\(\sigma_{\rm strength}/\rho\)、比弾性率は\(E/\rho\)である。対象の強度が引張・降伏のどちらかも示す。令和6年度基礎4-4はこの二つの取り違えを問う。

室温代表値では電気抵抗率はCuよりAu、AuよりFeが大きい\cite{openstax_resistivity}。純度と温度をそろえる。ステンレス鋼の一般的な定義はCr最低10.5質量\%、C最大1.2質量\%である\cite{stainless_definition2022}。したがってCr10\%未満は該当しないが、「10\%以上なら全て該当」という逆命題は導けない。

ダイヤモンドの高硬度を切削工具へ利用する例がある\cite{nims_diamond2007}。用途に応じて熱条件・被削材との化学反応も評価する。水素脆化は水素の侵入に伴う延性・破壊抵抗等の機械特性の劣化である。「物性が変化する」というだけでは範囲が広すぎ、改善や無関係な変化まで脆化とは呼ばない\cite{nims_hydrogen2024}。

硫黄快削鋼ではMnS等の介在物を利用して被削性を改善する。ただし工具寿命と仕上面粗さは同じ方向に良くなるとは限らず、脱酸・介在物の形・組織・切削速度も影響する\cite{sulfur_machinability2004}。S量だけから全鋼の被削性を保証せず、延性・溶接性など他の性能との交換条件を考える。

\section{水中の酸性度を共役塩基から考える}
\label{theory:basicchem:acid}
酸の強さは水中でプロトンを渡す平衡であり、分析濃度そのものではない。強酸の希薄溶液ではほぼ完全解離、弱酸では部分解離となる。単プロトン弱酸の分析濃度を\(c\)、解離定数を\(K_a\)とし、解離が十分小さければ\([\mathrm H^+]\simeq\sqrt{K_ac}\)である。水の自己解離が支配するほどの希薄条件や活量の影響を無視できない高濃度にはそのまま使わない。
共役塩基の安定化から構造の影響を読む。カルボキシラートの共鳴と電子求引基は解離を助け、CF3COOHは酢酸より強い。フェノキシドは共鳴安定化するためフェノールはエタノールより強い。HFとHClの比較では結合切断と溶媒和も入り、電気陰性度だけでは順位が決まらない\autocite{purdue_acid_strength,yale_hydrohalic}。炭酸の第一解離と溶存CO2を含む見かけの定義を分ける。

\section{生成エンタルピーの差として反応熱を計算する}
\label{theory:basicchem:hess}
エンタルピーは状態量なので、実際の反応経路によらず始状態と終状態の差で決まる。生成エンタルピーを用いると
\begin{equation}\Delta_r H^\circ=\sum_{\mathrm{products}}\nu_i\Delta_fH_i^\circ-\sum_{\mathrm{reactants}}\nu_i\Delta_fH_i^\circ.\end{equation}
係数と相を指定して足す。原典が生成時の放出熱を正の大きさで示す場合は、対応する生成エンタルピーを負に読み替える。液体水と水蒸気では値が違い、反応熱が小さい差で決まる設問では発熱・吸熱まで変わり得る\autocite{openstax_hess}。単位は指定反応式を一回進める基準か、生成物1 mol当たりかを明記する。

\section{沈殿した後に溶ける両性水酸化物}
\label{theory:basicchem:amphoteric}
Al3+にOH−を加えるとAl(OH)3が沈殿し、さらに過剰OH−でヒドロキソ錯イオンを作って溶ける。
\begin{equation}\mathrm{Al^{3+}+3OH^-\to Al(OH)_3},\qquad
\mathrm{Al(OH)_3+OH^-\rightleftharpoons[Al(OH)_4]^-}.\end{equation}
酸にも塩基にも反応する両性と、任意の沈殿が過剰試薬で溶けることは同じでない。添加剤がNaOHかアンモニアかを確認し、錯形成の条件を取り違えない\autocite{chem55_ampho}。

\section{熱伝導を担い手・散乱・空隙へ分ける}
\label{theory:basicchem:thermal}
通常の金属の熱輸送は自由電子の寄与が大きく、絶縁性結晶では格子振動が主要な担い手となる。合金化や構造の乱れは散乱を増す一般的傾向を持つ。低い熱伝導率の空隙を増せば通常は有効熱伝導率が低下する。ただし気孔中の対流・放射が重要な高温条件や配向の違いは別に検討する。
高分子の輸送も結晶化率だけでなく鎖の配向・界面・温度に依存する。結晶質と非晶質の通常比較を、任意の異組成間の順位へ広げない\autocite{aist_metal_thermal,openstax_thermal_conduction,chem55_polymer}。

\section{平衡収率と反応速度を分ける}
\label{theory:basicchem:equilibrium}
発熱のアンモニア合成\(\mathrm{N_2+3H_2\rightleftharpoons2NH_3}\)では、低温が生成側の平衡に有利である。体積を減らして圧力を上げると、気体4 molから2 molとなる生成側へ移る。温度は平衡定数を変え、圧縮は分圧と反応商を変える\autocite{chem55_equilibrium}。
触媒は平衡へ達する速度を変えるが平衡定数を変えない。低温による収率の増大と速度低下は別の効果なので、工業条件は両方と設備条件を考えて選ぶ。

\section{中和後に残るものからpHを判断する}
\label{theory:basicchem:neutralization}
体積と濃度から物質量を求め、放出・受容できるプロトンの当量を比較する。酸と塩基が当量でも、できた塩の溶液が必ず中性になるわけではない。弱塩基の共役酸NH4+は水へプロトンを渡して酸性を示す\autocite{chem55_hydrolysis}。
硫酸1 molとNaOH2 molは硫酸ナトリウムとなり、通常の中性近似又は硫酸イオンの弱い塩基性で扱う。硫酸1 molとNH3 2 molでは硫酸アンモニウムとなりNH4+の酸性が支配する。強塩基が過剰ならそのOH−を先に扱い、NH3とNH4+が共存する場合は緩衝系を考える。体積加算の近似は精密な濃度計算と区別する。

\section{腐食を環境と表面状態から判断する}
\label{theory:basicchem:corrosion}
腐食は材料と環境の反応で進む。イオン化傾向が大きくても、Alの強固な酸化膜のように反応を抑える不動態があれば通常の環境で耐食性を示す。ステンレスのCrによる不動態も同じ考え方だが、塩化物は孔食・すきま腐食を引き起こし得る\autocite{chem55_corrosion,stainless_definition2022}。海水用途は鋼種・温度・流動・形状等を合わせて選ぶ。
応力腐食割れは特定の材料・環境と引張応力の組合せで起こり、通常の引張強さより低い応力でも破壊し得る。水素脆化は水素侵入・拡散・欠陥との相互作用に関係する\autocite{nims_hydrogen2024}。これらを単なる均一減肉やすべての材料の性質へまとめない。

\section{金属資源の存在量と製錬反応}
\label{theory:basicchem:metallurgy}
地殻の存在量と工業生産量は違う。Alは地殻中で最も多い金属だが、金属生産量の順位をそこから直接決めない。反応性の高いAlは溶融塩中の酸化物を電解して得る\autocite{chem55_crust}。Tiの代表的なクロール法はTiCl4をMgで還元する\autocite{chem55_kroll}。
高炉では鉱石をコークス由来のCO等で還元し、石灰石等で不要成分を分離する\autocite{chem55_iron}。理想化した純粋Fe2O3から純粋Feを得る質量収支は、式量160のうちFeが112なので\(w_{\mathrm{ore}}=(160/112)w_{\mathrm{Fe}}\)。実操業では鉱石品位・歩留まり・銑鉄中の炭素を別に扱う。埋蔵量を生産量で割った可採年数は特定時点の指標であり、不変の枯渇期限ではない。

\section{密度・抵抗率・融点を独立に並べる}
\label{theory:basicchem:properties}
密度は単位体積当たり質量、抵抗率は形状を除いた電気輸送の指標、融点は指定圧力での相転移温度である。純金属の通常値では密度はCu\(>\)Fe\(>\)Al、20℃付近の抵抗率はFe\(>\)Al\(>\)Cu、融点はFe\(>\)Cu\(>\)Alとなる\autocite{openstax_resistivity,chem55_properties,chem55_properties_al,chem55_properties_fe}。三つの順位を一つの「金属の強さ」へ置き換えない。記号表の割当は年度ごとに戻して照合する。

\section{沸点上昇は溶質粒子数を数える}
\label{theory:basicchem:colligative}
不揮発性溶質を溶かすと蒸気圧が下がり、同じ外圧へ達するには温度を上げる必要がある。希薄な理想溶液の近似では\(\Delta T_b=K_b i b\)。\(b\)は溶媒1 kg当たりの溶質式単位の物質量、\(i\)は溶解後の有効粒子数の倍率である\autocite{chem55_colligative}。
同じ質量モル濃度なら、理想的な完全解離モデルの尿素等は1、NaClは2、CaCl2は3の粒子に対応する。実溶液の活量やイオン間相互作用により厳密な倍率からずれる。100℃を基準とする水の沸点は1気圧の条件付きであり、外圧が違えば純水の基準温度も変わる。

\section{材料名から代表的構成元素を読む}
\label{theory:basicchem:elements}
乾電池の代表的負極Zn、シリカ光ファイバーのSi、Al-Cu系ジュラルミンのCu、代表的な磁石のFeを用途と対応させる\autocite{chem55_fiber,chem55_duralumin,openstax_magnetism}。Liイオン電池ではCoを含む正極材料が代表例となるが、Coを含まない正極もあり、全電池共通の必要元素ではない\autocite{openstax_batteries}。材料の例と全称命題を区別し、正極と負極を読み分ける。

\section{ギブズエネルギーの単位をそろえる}
\label{theory:basicchem:gibbs}
\begin{equation}G=H-TS,\qquad\Delta_rG^\circ=\Delta_rH^\circ-T\Delta_rS^\circ.\end{equation}
\(T\)は絶対温度で、エントロピーの単位はエネルギー/温度である。したがって\(TS\)と\(H\)は同じ単位になる\autocite{chem55_gibbs}。一定温度・圧力での反応の駆動力を考えるとき、標準反応量と実際の組成での反応量を分け、後者には反応商の項が入る。

\section{結合の前後から有機反応を分類する}
\label{theory:basicchem:organic}
付加は多重結合に試薬が加わり、脱離は小分子が抜けて多重結合等ができ、置換は一つの原子団が別の原子団へ変わる\autocite{chem55_organic}。アルコールのOHからBrへの変換は置換、アルコールから水が抜ける反応は脱離、アルケンへのHBrは付加となる。カルボン酸からエステルへの変換はアシル基の置換として扱え、同時に縮合という観点でも呼べる。触媒名だけでなく、失われた結合と形成した結合を追う。

\section{反応の分類と反応が進む条件を分ける}
\label{theory:basicchem:redox}
酸化数の増加は酸化、減少は還元で、同一反応では両者を伴う。単体は0、通常の酸化物でOは−2、通常の酸や水でHは+1として電荷収支を取り、例外を確認する\autocite{iupac_oxidation_state}。
反応式の全元素で酸化数が不変ならその式は酸化還元ではない。ただし式が原子数と電荷を保存することは、その反応が指定条件で進むことの十分条件ではない。形式的な塩の交換式では溶解平衡や沈殿の駆動力を別に調べる。

\section{放射線の種類と材料への作用}
\label{theory:basicchem:radiation}
\(\alpha\)線はHe原子核、\(\beta^-\)線は電子、\(\beta^+\)線は陽電子、\(\gamma\)線は光子である。通常\(\alpha\)は強い電離作用と短い飛程を持つ。光子遮蔽の減衰は\(I=I_0e^{-\mu x}\)という単純モデルで表せ、\(\mu\)は材料と光子エネルギーに依存する\autocite{finish56_radiation}。同じ厚さと同じ面密度は別の比較である。照射は結晶欠陥や電子状態を変えて力学・電気特性を変え得る。

\section{構造と用途を対応させる化学物質}
\label{theory:basicchem:substances}
同素体は元素が同じでも結合・構造が違う。黒鉛は層状で軟らかく電気を伝え、ダイヤモンドの硬さ・通常の絶縁性と異なる。Haber--Boschはアンモニア合成、Solvayは炭酸ナトリウムの代表的製法である。
TiO2は光照射による表面反応を利用した抗菌・防汚等の光触媒材料として用いられる\autocite{finish56_titania}。性能は結晶相・複合化・光の波長等に依存する。鉛蓄電池の放電時は負極Pb、正極PbO2で、PbOという別酸化物と混同しない\autocite{openstax_batteries}。

\section{結合極性と分子極性を分ける}
\label{theory:basicchem:polarity}
極性結合を持つ分子でも、各結合の双極子が対称に配置されれば分子の双極子は0となる\autocite{finish56_polarity}。直線CO2、正四面体CH4・CCl4、平面三角形BF3は通常の形で相殺する。エーテルのO付近は折れ曲がり、分子の双極子が残る。結合を作る元素だけでなく、孤立電子対を含む分子形を確かめる。

\section{コロイドの粒子・界面・状態}
\label{theory:basicchem:colloids}
コロイドは分子溶液より大きい粒子等が分散する系で、粒子の界面と電荷が安定性を左右する。疎水コロイドに電解質を加えて反発を弱め、集合・沈殿させるのが凝析。小分子・イオンを膜で通して粒子と分けるのが透析である\autocite{openstax_colloids}。
光散乱で光路が見えるTyndall現象と、電場中で荷電粒子が移動する電気泳動は別の作用である。流動する分散系のゾルと、網目等ができ流動しにくいゲルを区別する。全コロイドが同じ量の電解質で凝析するわけではなく、表面・溶媒・イオン種等が関わる。

\subsection{平成23年度 問4-1}
\label{q:h23-reg-basic-4-1}

\noindent\textbf{原典：} sources/H23/basic.pdf、PDF物理12頁。

\textbf{公式正答：} 2（single）。

\par\textit{原典本文の校正と独自解説は未作成。分類のみの登録項目。}

\subsection{平成23年度 問4-2}
\label{q:h23-reg-basic-4-2}

\noindent\textbf{原典：} sources/H23/basic.pdf、PDF物理12頁。

\textbf{公式正答：} 1（single）。

\par\textit{原典本文の校正と独自解説は未作成。分類のみの登録項目。}

\subsection{平成23年度 問4-3}
\label{q:h23-reg-basic-4-3}

\noindent\textbf{原典：} sources/H23/basic.pdf、PDF物理12頁。

\textbf{公式正答：} 4（single）。

\par\textit{原典本文の校正と独自解説は未作成。分類のみの登録項目。}

\subsection{平成23年度 問4-4}
\label{q:h23-reg-basic-4-4}

\noindent\textbf{原典：} sources/H23/basic.pdf、PDF物理13頁。

\textbf{公式正答：} 4（single）。

\par\textit{原典本文の校正と独自解説は未作成。分類のみの登録項目。}

\subsection{平成23年度 問4-5}
\label{q:h23-reg-basic-4-5}

\noindent\textbf{原典：} sources/H23/basic.pdf、PDF物理14頁。

\textbf{公式正答：} 3（single）。

\par\textit{原典本文の校正と独自解説は未作成。分類のみの登録項目。}

\subsection{平成24年度 問4-1}
\label{q:h24-reg-basic-4-1}

\noindent\textbf{原典：} sources/H24/basic.pdf、PDF物理12頁。

\textbf{公式正答：} 1（single）。

\par\textit{原典本文の校正と独自解説は未作成。分類のみの登録項目。}

\subsection{平成24年度 問4-2}
\label{q:h24-reg-basic-4-2}

\noindent\textbf{原典：} sources/H24/basic.pdf、PDF物理12頁。

\textbf{公式正答：} 3（single）。

\par\textit{原典本文の校正と独自解説は未作成。分類のみの登録項目。}

\subsection{平成24年度 問4-3}
\label{q:h24-reg-basic-4-3}

\noindent\textbf{原典：} sources/H24/basic.pdf、PDF物理13頁。

\textbf{公式正答：} 2（single）。

\par\textit{原典本文の校正と独自解説は未作成。分類のみの登録項目。}

\subsection{平成24年度 問4-4}
\label{q:h24-reg-basic-4-4}

\noindent\textbf{原典：} sources/H24/basic.pdf、PDF物理13頁。

\textbf{公式正答：} 4（single）。

\par\textit{原典本文の校正と独自解説は未作成。分類のみの登録項目。}

\subsection{平成24年度 問4-5}
\label{q:h24-reg-basic-4-5}

\noindent\textbf{原典：} sources/H24/basic.pdf、PDF物理14頁。

\textbf{公式正答：} 3（single）。

\par\textit{原典本文の校正と独自解説は未作成。分類のみの登録項目。}

\subsection{平成25年度 問4-1}
\label{q:h25-reg-basic-4-1}

\noindent\textbf{原典：} sources/H25/basic.pdf、PDF物理15頁。

\textbf{公式正答：} 2（single）。

\par\textit{原典本文の校正と独自解説は未作成。分類のみの登録項目。}

\subsection{平成25年度 問4-2}
\label{q:h25-reg-basic-4-2}

\noindent\textbf{原典：} sources/H25/basic.pdf、PDF物理15頁。

\textbf{公式正答：} 2（single）。

\par\textit{原典本文の校正と独自解説は未作成。分類のみの登録項目。}

\subsection{平成25年度 問4-3}
\label{q:h25-reg-basic-4-3}

\noindent\textbf{原典：} sources/H25/basic.pdf、PDF物理16頁。

\textbf{公式正答：} 4（single）。

\par\textit{原典本文の校正と独自解説は未作成。分類のみの登録項目。}

\subsection{平成25年度 問4-4}
\label{q:h25-reg-basic-4-4}

\noindent\textbf{原典：} sources/H25/basic.pdf、PDF物理17頁。

\textbf{公式正答：} 3（single）。

\par\textit{原典本文の校正と独自解説は未作成。分類のみの登録項目。}

\subsection{平成25年度 問4-5}
\label{q:h25-reg-basic-4-5}

\noindent\textbf{原典：} sources/H25/basic.pdf、PDF物理17頁。

\textbf{公式正答：} 1（single）。

\par\textit{原典本文の校正と独自解説は未作成。分類のみの登録項目。}

\subsection{平成25年度 問4-6}
\label{q:h25-reg-basic-4-6}

\noindent\textbf{原典：} sources/H25/basic.pdf、PDF物理18頁。

\textbf{公式正答：} 3（single）。

\par\textit{原典本文の校正と独自解説は未作成。分類のみの登録項目。}

\subsection{平成26年度 問4-1}
\label{q:h26-reg-basic-4-1}

\noindent\textbf{原典：} sources/H26/basic.pdf、PDF物理13頁。

\textbf{公式正答：} 1（single）。

\par\textit{原典本文の校正と独自解説は未作成。分類のみの登録項目。}

\subsection{平成26年度 問4-2}
\label{q:h26-reg-basic-4-2}

\noindent\textbf{原典：} sources/H26/basic.pdf、PDF物理13頁。

\textbf{公式正答：} 2（single）。

\par\textit{原典本文の校正と独自解説は未作成。分類のみの登録項目。}

\subsection{平成26年度 問4-3}
\label{q:h26-reg-basic-4-3}

\noindent\textbf{原典：} sources/H26/basic.pdf、PDF物理14頁。

\textbf{公式正答：} 5（single）。

\par\textit{原典本文の校正と独自解説は未作成。分類のみの登録項目。}

\subsection{平成26年度 問4-4}
\label{q:h26-reg-basic-4-4}

\noindent\textbf{原典：} sources/H26/basic.pdf、PDF物理14頁。

\textbf{公式正答：} 4（single）。

\par\textit{原典本文の校正と独自解説は未作成。分類のみの登録項目。}

\subsection{平成26年度 問4-5}
\label{q:h26-reg-basic-4-5}

\noindent\textbf{原典：} sources/H26/basic.pdf、PDF物理15頁。

\textbf{公式正答：} 3（single）。

\par\textit{原典本文の校正と独自解説は未作成。分類のみの登録項目。}

\subsection{平成26年度 問4-6}
\label{q:h26-reg-basic-4-6}

\noindent\textbf{原典：} sources/H26/basic.pdf、PDF物理15頁。

\textbf{公式正答：} 1（single）。

\par\textit{原典本文の校正と独自解説は未作成。分類のみの登録項目。}

\subsection{平成27年度 問4-1}
\label{q:h27-reg-basic-4-1}

\noindent\textbf{原典：} sources/H27/basic.pdf、PDF物理12頁。

\textbf{公式正答：} 4（single）。

\par\textit{原典本文の校正と独自解説は未作成。分類のみの登録項目。}

\subsection{平成27年度 問4-2}
\label{q:h27-reg-basic-4-2}

\noindent\textbf{原典：} sources/H27/basic.pdf、PDF物理12頁。

\textbf{公式正答：} 3（single）。

\par\textit{原典本文の校正と独自解説は未作成。分類のみの登録項目。}

\subsection{平成27年度 問4-3}
\label{q:h27-reg-basic-4-3}

\noindent\textbf{原典：} sources/H27/basic.pdf、PDF物理12頁。

\textbf{公式正答：} 5（single）。

\par\textit{原典本文の校正と独自解説は未作成。分類のみの登録項目。}

\subsection{平成27年度 問4-4}
\label{q:h27-reg-basic-4-4}

\noindent\textbf{原典：} sources/H27/basic.pdf、PDF物理13頁。

\textbf{公式正答：} 1（single）。

\par\textit{原典本文の校正と独自解説は未作成。分類のみの登録項目。}

\subsection{平成27年度 問4-5}
\label{q:h27-reg-basic-4-5}

\noindent\textbf{原典：} sources/H27/basic.pdf、PDF物理14頁。

\textbf{公式正答：} 2（single）。

\par\textit{原典本文の校正と独自解説は未作成。分類のみの登録項目。}

\subsection{平成27年度 問4-6}
\label{q:h27-reg-basic-4-6}

\noindent\textbf{原典：} sources/H27/basic.pdf、PDF物理14頁。

\textbf{公式正答：} 3（single）。

\par\textit{原典本文の校正と独自解説は未作成。分類のみの登録項目。}

\subsection{平成28年度 問4-1}
\label{q:h28-reg-basic-4-1}

\noindent\textbf{原典：} sources/H28/basic.pdf、PDF物理17頁。

\textbf{公式正答：} 4（single）。

\par\textit{原典本文の校正と独自解説は未作成。分類のみの登録項目。}

\subsection{平成28年度 問4-2}
\label{q:h28-reg-basic-4-2}

\noindent\textbf{原典：} sources/H28/basic.pdf、PDF物理17頁。

\textbf{公式正答：} 4（single）。

\par\textit{原典本文の校正と独自解説は未作成。分類のみの登録項目。}

\subsection{平成28年度 問4-3}
\label{q:h28-reg-basic-4-3}

\noindent\textbf{原典：} sources/H28/basic.pdf、PDF物理18頁。

\textbf{公式正答：} 3（single）。

\par\textit{原典本文の校正と独自解説は未作成。分類のみの登録項目。}

\subsection{平成28年度 問4-4}
\label{q:h28-reg-basic-4-4}

\noindent\textbf{原典：} sources/H28/basic.pdf、PDF物理18頁。

\textbf{公式正答：} 2（single）。

\par\textit{原典本文の校正と独自解説は未作成。分類のみの登録項目。}

\subsection{平成28年度 問4-5}
\label{q:h28-reg-basic-4-5}

\noindent\textbf{原典：} sources/H28/basic.pdf、PDF物理19頁。

\textbf{公式正答：} 1（single）。

\par\textit{原典本文の校正と独自解説は未作成。分類のみの登録項目。}

\subsection{平成28年度 問4-6}
\label{q:h28-reg-basic-4-6}

\noindent\textbf{原典：} sources/H28/basic.pdf、PDF物理20頁。

\textbf{公式正答：} 3（single）。

\par\textit{原典本文の校正と独自解説は未作成。分類のみの登録項目。}

\subsection{平成29年度 問4-1}
\label{q:h29-reg-basic-4-1}

\noindent\textbf{原典：} sources/H29/basic.pdf、PDF物理15頁。

\textbf{公式正答：} 4（single）。

\par\textit{原典本文の校正と独自解説は未作成。分類のみの登録項目。}

\subsection{平成29年度 問4-2}
\label{q:h29-reg-basic-4-2}

\noindent\textbf{原典：} sources/H29/basic.pdf、PDF物理15頁。

\textbf{公式正答：} 5（single）。

\par\textit{原典本文の校正と独自解説は未作成。分類のみの登録項目。}

\subsection{平成29年度 問4-3}
\label{q:h29-reg-basic-4-3}

\noindent\textbf{原典：} sources/H29/basic.pdf、PDF物理16頁。

\textbf{公式正答：} 4（single）。

\par\textit{原典本文の校正と独自解説は未作成。分類のみの登録項目。}

\subsection{平成29年度 問4-4}
\label{q:h29-reg-basic-4-4}

\noindent\textbf{原典：} sources/H29/basic.pdf、PDF物理16頁。

\textbf{公式正答：} 1（single）。

\par\textit{原典本文の校正と独自解説は未作成。分類のみの登録項目。}

\subsection{平成29年度 問4-5}
\label{q:h29-reg-basic-4-5}

\noindent\textbf{原典：} sources/H29/basic.pdf、PDF物理17頁。

\textbf{公式正答：} 1（single）。

\par\textit{原典本文の校正と独自解説は未作成。分類のみの登録項目。}

\subsection{平成29年度 問4-6}
\label{q:h29-reg-basic-4-6}

\noindent\textbf{原典：} sources/H29/basic.pdf、PDF物理17頁。

\textbf{公式正答：} 1（single）。

\par\textit{原典本文の校正と独自解説は未作成。分類のみの登録項目。}

\subsection{平成30年度 問4-1}
\label{q:h30-reg-basic-4-1}

\noindent\textbf{原典：} sources/H30/basic.pdf、PDF物理15頁。

\textbf{公式正答：} 5（single）。

\par\textit{原典本文の校正と独自解説は未作成。分類のみの登録項目。}

\subsection{平成30年度 問4-2}
\label{q:h30-reg-basic-4-2}

\noindent\textbf{原典：} sources/H30/basic.pdf、PDF物理15頁。

\textbf{公式正答：} 5（single）。

\par\textit{原典本文の校正と独自解説は未作成。分類のみの登録項目。}

\subsection{平成30年度 問4-3}
\label{q:h30-reg-basic-4-3}

\noindent\textbf{原典：} sources/H30/basic.pdf、PDF物理15頁。

\textbf{公式正答：} 3（single）。

\par\textit{原典本文の校正と独自解説は未作成。分類のみの登録項目。}

\subsection{平成30年度 問4-4}
\label{q:h30-reg-basic-4-4}

\noindent\textbf{原典：} sources/H30/basic.pdf、PDF物理16頁。

\textbf{公式正答：} 2（single）。

\par\textit{原典本文の校正と独自解説は未作成。分類のみの登録項目。}

\subsection{平成30年度 問4-5}
\label{q:h30-reg-basic-4-5}

\noindent\textbf{原典：} sources/H30/basic.pdf、PDF物理16頁。

\textbf{公式正答：} 5（single）。

\par\textit{原典本文の校正と独自解説は未作成。分類のみの登録項目。}

\subsection{平成30年度 問4-6}
\label{q:h30-reg-basic-4-6}

\noindent\textbf{原典：} sources/H30/basic.pdf、PDF物理17頁。

\textbf{公式正答：} 3（single）。

\par\textit{原典本文の校正と独自解説は未作成。分類のみの登録項目。}

\subsection{令和元年度 問4-1}
\label{q:r01-reg-basic-4-1}

\noindent\textbf{原典：} sources/R01/basic.pdf、PDF物理16頁。

\textbf{公式正答：} 4（single）。

\par\textit{原典本文の校正と独自解説は未作成。分類のみの登録項目。}

\subsection{令和元年度 問4-2}
\label{q:r01-reg-basic-4-2}

\noindent\textbf{原典：} sources/R01/basic.pdf、PDF物理16頁。

\textbf{公式正答：} 5（single）。

\par\textit{原典本文の校正と独自解説は未作成。分類のみの登録項目。}

\subsection{令和元年度 問4-3}
\label{q:r01-reg-basic-4-3}

\noindent\textbf{原典：} sources/R01/basic.pdf、PDF物理17頁。

\textbf{公式正答：} 4（single）。

\par\textit{原典本文の校正と独自解説は未作成。分類のみの登録項目。}

\subsection{令和元年度 問4-4}
\label{q:r01-reg-basic-4-4}

\noindent\textbf{原典：} sources/R01/basic.pdf、PDF物理17頁。

\textbf{公式正答：} 3（single）。

\par\textit{原典本文の校正と独自解説は未作成。分類のみの登録項目。}

\subsection{令和元年度 問4-5}
\label{q:r01-reg-basic-4-5}

\noindent\textbf{原典：} sources/R01/basic.pdf、PDF物理18頁。

\textbf{公式正答：} 4（single）。

\par\textit{原典本文の校正と独自解説は未作成。分類のみの登録項目。}

\subsection{令和元年度 問4-6}
\label{q:r01-reg-basic-4-6}

\noindent\textbf{原典：} sources/R01/basic.pdf、PDF物理18頁。

\textbf{公式正答：} 3（single）。

\par\textit{原典本文の校正と独自解説は未作成。分類のみの登録項目。}

\subsection{令和元年度再試験 問4-1}
\label{q:r01-ret-basic-4-1}

\noindent\textbf{原典：} sources/R01\_retest/basic.pdf、PDF物理14頁。

\textbf{公式正答：} 2（single）。

\par\textit{原典本文の校正と独自解説は未作成。分類のみの登録項目。}

\subsection{令和元年度再試験 問4-2}
\label{q:r01-ret-basic-4-2}

\noindent\textbf{原典：} sources/R01\_retest/basic.pdf、PDF物理14頁。

\textbf{公式正答：} 3（single）。

\par\textit{原典本文の校正と独自解説は未作成。分類のみの登録項目。}

\subsection{令和元年度再試験 問4-3}
\label{q:r01-ret-basic-4-3}

\noindent\textbf{原典：} sources/R01\_retest/basic.pdf、PDF物理14頁。

\textbf{公式正答：} 2（single）。

\par\textit{原典本文の校正と独自解説は未作成。分類のみの登録項目。}

\subsection{令和元年度再試験 問4-4}
\label{q:r01-ret-basic-4-4}

\noindent\textbf{原典：} sources/R01\_retest/basic.pdf、PDF物理15頁。

\textbf{公式正答：} 4（single）。

\par\textit{原典本文の校正と独自解説は未作成。分類のみの登録項目。}

\subsection{令和元年度再試験 問4-5}
\label{q:r01-ret-basic-4-5}

\noindent\textbf{原典：} sources/R01\_retest/basic.pdf、PDF物理15頁。

\textbf{公式正答：} 3（single）。

\par\textit{原典本文の校正と独自解説は未作成。分類のみの登録項目。}

\subsection{令和元年度再試験 問4-6}
\label{q:r01-ret-basic-4-6}

\noindent\textbf{原典：} sources/R01\_retest/basic.pdf、PDF物理16頁。

\textbf{公式正答：} 3（single）。

\par\textit{原典本文の校正と独自解説は未作成。分類のみの登録項目。}

\subsection{令和2年度 問4-1}
\label{q:r02-reg-basic-4-1}

\noindent\textbf{原典：} sources/R02/basic.pdf、PDF物理15頁。

\textbf{公式正答：} 2（single）。

\par\textit{原典本文の校正と独自解説は未作成。分類のみの登録項目。}

\subsection{令和2年度 問4-2}
\label{q:r02-reg-basic-4-2}

\noindent\textbf{原典：} sources/R02/basic.pdf、PDF物理15頁。

\textbf{公式正答：} 3（single）。

\par\textit{原典本文の校正と独自解説は未作成。分類のみの登録項目。}

\subsection{令和2年度 問4-3}
\label{q:r02-reg-basic-4-3}

\noindent\textbf{原典：} sources/R02/basic.pdf、PDF物理16頁。

\textbf{公式正答：} 4（single）。

\par\textit{原典本文の校正と独自解説は未作成。分類のみの登録項目。}

\subsection{令和2年度 問4-4}
\label{q:r02-reg-basic-4-4}

\noindent\textbf{原典：} sources/R02/basic.pdf、PDF物理17頁。

\textbf{公式正答：} 5（single）。

\par\textit{原典本文の校正と独自解説は未作成。分類のみの登録項目。}

\subsection{令和2年度 問4-5}
\label{q:r02-reg-basic-4-5}

\noindent\textbf{原典：} sources/R02/basic.pdf、PDF物理18頁。

\textbf{公式正答：} 3（single）。

\par\textit{原典本文の校正と独自解説は未作成。分類のみの登録項目。}

\subsection{令和2年度 問4-6}
\label{q:r02-reg-basic-4-6}

\noindent\textbf{原典：} sources/R02/basic.pdf、PDF物理18頁。

\textbf{公式正答：} 4（single）。

\par\textit{原典本文の校正と独自解説は未作成。分類のみの登録項目。}

\subsection{令和3年度 問4-1}
\label{q:r03-reg-basic-4-1}

\noindent\textbf{原典：} sources/R03/basic.pdf、PDF物理16頁。

\textbf{公式正答：} 3（single）。

\par\textit{原典本文の校正と独自解説は未作成。分類のみの登録項目。}

\subsection{令和3年度 問4-2}
\label{q:r03-reg-basic-4-2}

\noindent\textbf{原典：} sources/R03/basic.pdf、PDF物理16頁。

\textbf{公式正答：} 4（single）。

\par\textit{原典本文の校正と独自解説は未作成。分類のみの登録項目。}

\subsection{令和3年度 問4-3}
\label{q:r03-reg-basic-4-3}

\noindent\textbf{原典：} sources/R03/basic.pdf、PDF物理17頁。

\textbf{公式正答：} 2（single）。

\par\textit{原典本文の校正と独自解説は未作成。分類のみの登録項目。}

\subsection{令和3年度 問4-4}
\label{q:r03-reg-basic-4-4}

\noindent\textbf{原典：} sources/R03/basic.pdf、PDF物理18頁。

\textbf{公式正答：} 1（single）。

\par\textit{原典本文の校正と独自解説は未作成。分類のみの登録項目。}

\subsection{令和3年度 問4-5}
\label{q:r03-reg-basic-4-5}

\noindent\textbf{原典：} sources/R03/basic.pdf、PDF物理19頁。

\textbf{公式正答：} 1（single）。

\par\textit{原典本文の校正と独自解説は未作成。分類のみの登録項目。}

\subsection{令和3年度 問4-6}
\label{q:r03-reg-basic-4-6}

\noindent\textbf{原典：} sources/R03/basic.pdf、PDF物理20頁。

\textbf{公式正答：} 2（single）。

\par\textit{原典本文の校正と独自解説は未作成。分類のみの登録項目。}

\subsection{令和4年度 問4-1}
\label{q:r04-reg-basic-4-1}

\noindent\textbf{原典：} sources/R04/basic.pdf、PDF物理14頁。

\textbf{公式正答：} 2（single）。

\par\textit{原典本文の校正と独自解説は未作成。分類のみの登録項目。}

\subsection{令和4年度 問4-2}
\label{q:r04-reg-basic-4-2}

\noindent\textbf{原典：} sources/R04/basic.pdf、PDF物理14頁。

\textbf{公式正答：} 4（single）。

\par\textit{原典本文の校正と独自解説は未作成。分類のみの登録項目。}

\subsection{令和4年度 問4-3}
\label{q:r04-reg-basic-4-3}

\noindent\textbf{原典：} sources/R04/basic.pdf、PDF物理14頁。

\textbf{公式正答：} 1（single）。

\par\textit{原典本文の校正と独自解説は未作成。分類のみの登録項目。}

\subsection{令和4年度 問4-4}
\label{q:r04-reg-basic-4-4}

\noindent\textbf{原典：} sources/R04/basic.pdf、PDF物理15頁。

\textbf{公式正答：} 1（single）。

\par\textit{原典本文の校正と独自解説は未作成。分類のみの登録項目。}

\subsection{令和4年度 問4-5}
\label{q:r04-reg-basic-4-5}

\noindent\textbf{原典：} sources/R04/basic.pdf、PDF物理15頁。

\textbf{公式正答：} 5（single）。

\par\textit{原典本文の校正と独自解説は未作成。分類のみの登録項目。}

\subsection{令和4年度 問4-6}
\label{q:r04-reg-basic-4-6}

\noindent\textbf{原典：} sources/R04/basic.pdf、PDF物理16頁。

\textbf{公式正答：} 5（single）。

\par\textit{原典本文の校正と独自解説は未作成。分類のみの登録項目。}

\subsection{令和5年度 問4-1}
\label{q:r05-reg-basic-4-1}

\noindent\textbf{原典：} sources/R05/basic.pdf、PDF物理13頁。

\textbf{公式正答：} 3（single）。

\par\textit{原典本文の校正と独自解説は未作成。分類のみの登録項目。}

\subsection{令和5年度 問4-2}
\label{q:r05-reg-basic-4-2}

\noindent\textbf{原典：} sources/R05/basic.pdf、PDF物理13頁。

\textbf{公式正答：} 5（single）。

\par\textit{原典本文の校正と独自解説は未作成。分類のみの登録項目。}

\subsection{令和5年度 問4-3}
\label{q:r05-reg-basic-4-3}

\noindent\textbf{原典：} sources/R05/basic.pdf、PDF物理13頁。

\textbf{公式正答：} 3（single）。

\par\textit{原典本文の校正と独自解説は未作成。分類のみの登録項目。}

\subsection{令和5年度 問4-4}
\label{q:r05-reg-basic-4-4}

\noindent\textbf{原典：} sources/R05/basic.pdf、PDF物理14頁。

\textbf{公式正答：} 1（single）。

\par\textit{原典本文の校正と独自解説は未作成。分類のみの登録項目。}

\subsection{令和5年度 問4-5}
\label{q:r05-reg-basic-4-5}

\noindent\textbf{原典：} sources/R05/basic.pdf、PDF物理14頁。

\textbf{公式正答：} 3（single）。

\par\textit{原典本文の校正と独自解説は未作成。分類のみの登録項目。}

\subsection{令和5年度 問4-6}
\label{q:r05-reg-basic-4-6}

\noindent\textbf{原典：} sources/R05/basic.pdf、PDF物理15頁。

\textbf{公式正答：} 5（single）。

\par\textit{原典本文の校正と独自解説は未作成。分類のみの登録項目。}

\subsection{令和6年度 問4-1}
\label{q:r06-reg-basic-4-1}

\noindent\textbf{原典：} sources/R06/basic.pdf、PDF物理15頁。

\textbf{公式正答：} 1（single）。

\par\textit{原典本文の校正と独自解説は未作成。分類のみの登録項目。}

\subsection{令和6年度 問4-2}
\label{q:r06-reg-basic-4-2}

\noindent\textbf{原典：} sources/R06/basic.pdf、PDF物理15頁。

\textbf{公式正答：} 3（single）。

\par\textit{原典本文の校正と独自解説は未作成。分類のみの登録項目。}

\subsection{令和6年度 問4-3}
\label{q:r06-reg-basic-4-3}

\noindent\textbf{原典：} sources/R06/basic.pdf、PDF物理15頁。

\textbf{公式正答：} 3（single）。

\par\textit{原典本文の校正と独自解説は未作成。分類のみの登録項目。}

\subsection{令和6年度 問4-4}
\label{q:r06-reg-basic-4-4}

\noindent\textbf{原典：} sources/R06/basic.pdf、PDF物理16頁。

\textbf{公式正答：} 2（single）。

\par\textit{原典本文の校正と独自解説は未作成。分類のみの登録項目。}

\subsection{令和6年度 問4-5}
\label{q:r06-reg-basic-4-5}

\noindent\textbf{原典：} sources/R06/basic.pdf、PDF物理16頁。

\textbf{公式正答：} 1（single）。

\par\textit{原典本文の校正と独自解説は未作成。分類のみの登録項目。}

\subsection{令和6年度 問4-6}
\label{q:r06-reg-basic-4-6}

\noindent\textbf{原典：} sources/R06/basic.pdf、PDF物理16頁。

\textbf{公式正答：} 4（single）。

\par\textit{原典本文の校正と独自解説は未作成。分類のみの登録項目。}

\subsection{令和7年度 問4-1}
\label{q:r07-reg-basic-4-1}

\noindent\textbf{原典：} sources/R07/basic.pdf、PDF物理13頁。

\textbf{公式正答：} 2（single）。

\par\textit{原典本文の校正と独自解説は未作成。分類のみの登録項目。}

\subsection{令和7年度 問4-2}
\label{q:r07-reg-basic-4-2}

\noindent\textbf{原典：} sources/R07/basic.pdf、PDF物理13頁。

\textbf{公式正答：} 3（single）。

\par\textit{原典本文の校正と独自解説は未作成。分類のみの登録項目。}

\subsection{令和7年度 問4-3}
\label{q:r07-reg-basic-4-3}

\noindent\textbf{原典：} sources/R07/basic.pdf、PDF物理13頁。

\textbf{公式正答：} 3（single）。

\par\textit{原典本文の校正と独自解説は未作成。分類のみの登録項目。}

\subsection{令和7年度 問4-4}
\label{q:r07-reg-basic-4-4}

\noindent\textbf{原典：} sources/R07/basic.pdf、PDF物理14頁。

\textbf{公式正答：} 5（single）。

\par\textit{原典本文の校正と独自解説は未作成。分類のみの登録項目。}

\subsection{令和7年度 問4-5}
\label{q:r07-reg-basic-4-5}

\noindent\textbf{原典：} sources/R07/basic.pdf、PDF物理14頁。

\textbf{公式正答：} 4（single）。

\par\textit{原典本文の校正と独自解説は未作成。分類のみの登録項目。}

\subsection{令和7年度 問4-6}
\label{q:r07-reg-basic-4-6}

\noindent\textbf{原典：} sources/R07/basic.pdf、PDF物理14頁。

\textbf{公式正答：} 4（single）。

\par\textit{原典本文の校正と独自解説は未作成。分類のみの登録項目。}

